\section{Conclusion}\label{sec:conclusion}
%
In the context of proof of space consensus, Signum's advantage is its simplicity,
the small size of its proofs (deadlines, around $200$ bytes)
and the absence of an initialization phase and transactions.
Moreover, it is the only fully implemented and deployed solution and
supports smart contracts.
Signum has recently added a layer of proof of commitment, that amplifies
the fees for miners that staked more crypto. This seems just a way to sustain
the value of their cyprotcurrency and is irrelevant \wrt the
underlying proof of space algorithm described in this paper.
There are, of course, drawbacks in Signum's algorithm.
First of all, it is
theoretically possible to mine new blocks in a proof of work style, although
this is not been observed in practice up to now.
More generally, there are drawbacks of proof of space that still need to
be investigated: waste of storage, SSD degradation and even storage farm
centralization.

Our formalization of Signum's consensus is valuable because it sheds light
on a blockchain network that runs since ten years but was missing any formal definition.
Moreover, it allowed us to understand that Signum is free from block-grinding attacks
and is largely protected from challenge-grinding attacks.
%This is currently not convenient since Alg.~\ref{alg:nonce_construction} is
%relatively slow, in particular
%by using the expensive shabal256 hashing for $h_\deadline$ (Tab.~\ref{tab:notations}).
%But ASICs might change the situation in the future.
%At the end, it will be the relative
%increase of ASICs speed and memory size that will decide if Signum remains a
%proof of space network or if it becomes more rational to mine it with proof of work.
